\documentclass[11pt]{article}

\usepackage[margin=1in]{geometry}
\usepackage{amsmath,amssymb,amsthm}
\usepackage{booktabs}
\usepackage{array}
\usepackage{enumitem}
\usepackage[colorlinks=true,linkcolor=blue,citecolor=blue,urlcolor=blue]{hyperref}

\setlist[itemize]{leftmargin=1.35em,itemsep=0.25em,topsep=0.25em}
\setlist[enumerate]{leftmargin=1.45em,itemsep=0.25em,topsep=0.25em}
\emergencystretch=3em

\newcommand{\Tset}{\mathcal T}
\newcommand{\CB}{\mathcal C_B}
\newcommand{\E}{\mathbb E}
\newcommand{\argmax}{\operatorname*{arg\,max}}

\newtheorem{proposition}{Proposition}
\newtheorem{lemma}{Lemma}

\title{\vspace{-1.5em}Research Direction Memo\\
\large When to Correct? Statistical Timing of Corrective Refinement}
\author{Supervisor-facing draft}
\date{11 May 2026}

\begin{document}
\maketitle

\vspace{-1.0em}
\section*{Proposed thesis question}

The thesis studies correction timing in masked diffusion as a finite-budget
statistical decision problem.  The predictor schedule, corrector mechanism,
and quality metric are fixed.  The object of study is when a limited number of
corrective refinement steps should be spent along the generation trajectory.

\medskip
\noindent
The main research question is:

\begin{quote}
When should a masked diffusion sampler spend a limited correction budget, and
can correction timing be predicted from the current generation state?
\end{quote}

\noindent
A more formal version is:

\begin{quote}
For masked diffusion language models with a finite number of corrective
refinement steps, can the time steps at which correction is most valuable be
characterized and predicted using observable uncertainty and error signals
from the generation trajectory?
\end{quote}

\noindent
The interaction question remains part of the thesis, but as an explanatory
sub-question rather than the headline:

\begin{quote}
Are correction gains approximately additive over time, or do redundancy and
interactions between correction steps explain why simple timing rules fail?
\end{quote}

\section*{Why this is the right framing}

This framing keeps the thesis close to statistics.  The core quantities are
potential outcomes, conditional expectations, held-out prediction error,
finite-budget regret certificates, and interaction diagnostics.  Search
algorithms are useful, but they are not the main contribution.  They act as
empirical probes showing whether useful correction-timing structure exists
beyond simple uncertainty rankers.

It also separates three questions that should not be conflated:
\begin{enumerate}
\item \textbf{Correction value:} where in generation does correction help?
\item \textbf{State predictability:} can observable trajectory statistics
predict that value beyond absolute time?
\item \textbf{Interaction structure:} if simple marginal rules fail, is the
failure caused by redundancy, higher-order interactions, or lack of stable
local structure?
\end{enumerate}

This is different from training a new corrector or studying which tokens to
revise.  PRISM, locally balanced proposals, and informed-corrector kernels are
important related work for how to design informed correction moves.  Here the
correction move is fixed; the thesis studies allocation of that move across
time.

\section*{Statistical setup}

Let \(\Tset=\{1,\ldots,T\}\).  A correction schedule is a subset
\(S\subseteq\Tset\), with \(|S|=B\).  Fix a masked-diffusion predictor, a
corrector kernel, a coupling variable \(\omega\) for common random numbers, and
an integrable terminal utility \(F:\mathcal Y\to\mathbb R\).  The final sample
under schedule \(S\) is
\[
Y^S=\Phi(S,\omega),
\]
and the population schedule value is
\[
J(S)=\E[F(Y^S)], \qquad G(S)=J(S)-J(\emptyset).
\]

For paired experiments, seed-level potential outcomes are
\[
G_i(S)=F(Y_i^S)-F(Y_i^\emptyset),
\]
where the same seed \(i\) is reused across schedules.  This makes correction
timing a potential-outcome problem: the counterfactual is what the same
generation trajectory would have produced with a correction inserted at a
different time.

\paragraph{Marginal correction value.}
For a single correction at time \(t\),
\[
\Delta_t=J(\{t\})-J(\emptyset), \qquad
\delta_{i,t}=G_i(\{t\}).
\]
\(\Delta_t\) answers the descriptive question: at which generation stages is a
single correction most valuable on average?

\paragraph{State-predictive correction value.}
Let \(Z_{i,t}\) be observable state features before deciding whether to
correct at time \(t\): diffusion phase, mask ratio, entropy, top-1/top-2
margin, recently unmasked fraction, token instability, likelihood change, or
pre/post prediction disagreement when available.  The statistical target is
\[
m(t,z)=\E[\delta_t \mid t, Z_t=z].
\]
The key question is whether \(Z_t\) improves prediction of \(\delta_t\) beyond
time \(t\) alone.

\paragraph{Budgeted timing rule.}
Given a predictor \(\widehat m(t,Z_t)\), the simplest offline timing rule
selects the top \(B\) time steps by predicted correction value.  An online
trigger is the sequential version: spend correction budget only when
\(\widehat m(t,Z_t)\) is high enough relative to the remaining time and
remaining budget.

\section*{Theorem candidates}

\begin{proposition}[State information as prediction gain]
Let \(\delta\in L^2\) be the marginal correction value for a randomly sampled
trajectory-time pair.  Let \(\mathcal F_0=\sigma(t)\) be the information in
time alone and let \(\mathcal F_1=\sigma(t,Z_t)\), with
\(\mathcal F_0\subseteq\mathcal F_1\).  The Bayes predictors under squared loss
are
\[
m_0(t)=\E[\delta\mid\mathcal F_0],
\qquad
m_1(t,Z_t)=\E[\delta\mid\mathcal F_1].
\]
Their mean-squared prediction errors satisfy
\[
\E[(\delta-m_0)^2]-\E[(\delta-m_1)^2]
= \E[(m_1-m_0)^2]\ge 0.
\]
\end{proposition}

\begin{proof}[Proof sketch]
Conditional expectation is the \(L^2\) projection onto the corresponding
information set.  Since \(\mathcal F_0\subseteq\mathcal F_1\), the
Pythagorean identity for nested projections gives the displayed equality.
\end{proof}

\noindent
\textbf{Use in the thesis.}  This gives the clean statistical meaning of the
state-predictability question.  If held-out prediction using \(Z_t\) does not
beat time-only prediction, then a state-adaptive trigger is not justified.  If
it does, the improvement measures how much information the current generation
state carries about correction value.

\begin{proposition}[Finite-pool regret certificate for predicted marginal timing]
Let \(\CB\subseteq\{S\subseteq\Tset: |S|=B\}\) be a fixed candidate pool and
let
\[
A(S)=\sum_{t\in S}\Delta_t,\qquad
\widehat A(S)=\sum_{t\in S}\widehat m(t,Z_t).
\]
Assume
\[
\sup_{S\in\CB}|G(S)-A(S)|\le \eta,
\qquad
\sup_{S\in\CB}|A(S)-\widehat A(S)|\le \alpha.
\]
If \(S_G^\star\in\argmax_{S\in\CB}G(S)\) and
\(\widehat S\in\argmax_{S\in\CB}\widehat A(S)\), then
\[
G(S_G^\star)-G(\widehat S)\le 2\eta+2\alpha .
\]
\end{proposition}

\begin{proof}[Proof sketch]
Apply the generic finite-pool surrogate argument with surrogate \(A\) and
estimated surrogate \(\widehat A\).  The factor \(2\) appears because the
surrogate error is paid once at the oracle schedule and once at the selected
schedule.
\end{proof}

\noindent
\textbf{Use in the thesis.}  This theorem clarifies when a lightweight timing
rule can be defended: correction gains must be approximately additive on the
candidate schedules and the state predictor must be accurate enough on
held-out trajectories.  If either condition fails, the failure is a
statistical result, not an implementation failure.

\begin{lemma}[Interaction diagnostic]
For \(s\neq t\), define the pair interaction
\[
\xi_{s,t}=G(\{s,t\})-\Delta_s-\Delta_t .
\]
For a quadratic approximation
\[
Q_2(S)=\sum_{t\in S}a_t+\sum_{\{s,t\}\subseteq S}b_{s,t},
\]
the diminishing-return gap for \(A\subseteq B\subseteq\Tset\) and
\(x\notin B\) is
\[
\bigl[Q_2(A\cup\{x\})-Q_2(A)\bigr]
-
\bigl[Q_2(B\cup\{x\})-Q_2(B)\bigr]
=-\sum_{t\in B\setminus A}b_{x,t}.
\]
Thus negative pair coefficients imply diminishing returns for the quadratic
approximation, but not automatically for the true \(G\).
\end{lemma}

\begin{proof}[Proof sketch]
The marginal gain of adding \(x\) to \(A\) is
\(a_x+\sum_{t\in A}b_{x,t}\).  Subtract the corresponding expression for
\(B\).  A claim about \(G\) needs residual control between \(G\) and \(Q_2\).
\end{proof}

\noindent
\textbf{Use in the thesis.}  This is the formal bridge from state-predicted
marginal timing to interaction diagnostics.  If marginal or state-adaptive
timing fails, \(\xi_{s,t}\), additivity residuals, and one-swap neighborhoods
explain whether the failure is due to redundancy, higher-order effects, or an
essentially black-box landscape.

\section*{Empirical plan}

\paragraph{RQ1: Where is correction valuable?}
Estimate \(\Delta_t\) over generation time using paired common-random-number
runs.  Report curves, confidence intervals, and phase summaries.

\paragraph{RQ2: Is correction value predictable from state?}
Compare time-only prediction \(m_0(t)\) against state-based prediction
\(m_1(t,Z_t)\).  Features can include mask ratio, entropy, margin, recently
unmasked fraction, token instability, likelihood change, and prediction
disagreement.  Use held-out trajectories; do not evaluate a trigger on the
same data used to choose it.

\paragraph{RQ3: Can simple timing rules spend budget well?}
Compare uniform, early-heavy, late-heavy, marginal-greedy,
state-triggered, and oracle schedules under the same budget \(B\).  A learned
logistic or small MLP trigger is optional and only justified if RQ2 shows
held-out predictive signal.

\paragraph{RQ4: Why do simple rules fail?}
Measure additivity residuals, pair interactions, diminishing-return gaps, and
one-swap neighborhood improvements.  This is the interaction/redundancy part
of the thesis.

\section*{What current evidence already answers}

Many of the empirical questions are already partly answered on ProSeCo-OWT.
The thesis does not need to pretend these are open; the supervisor-facing
story can present the current evidence as the beginning of the result.

\begin{center}
\small
\begin{tabular}{p{0.20\linewidth}p{0.51\linewidth}p{0.21\linewidth}}
\toprule
Question & Current answer on ProSeCo-OWT & Status \\
\midrule
Is correction timing non-vacuous? &
Yes.  Existing paired experiments show MC-oracle headroom over uniform of
about \(+0.45\) at \(B\in\{2,3,4\}\), with a K=30 replication around \(+0.37\). &
Answered for ProSeCo-OWT. \\
\addlinespace
Do simple uncertainty/ranker signals recover the headroom? &
No.  Tested separable rankers fail to recover the MC-oracle headroom.  The
mean marginal envelope becomes weak by larger budgets. &
Answered negatively for tested signals. \\
\addlinespace
Can simple adaptive state conditioning solve it? &
Current evidence says no.  The bucketed adaptive-controller pilot shrinks
calibration error by less than \(1.7\%\) and does not recover headroom.  The
held-out state-predictability audit also finds no improvement from adding
current state features to time-only prediction. &
Answered negatively for current features. \\
\addlinespace
Is there useful structure beyond simple rankers? &
Yes.  True-\(G\) search procedures recover much of the oracle headroom:
CD-G about \(74\)--\(84\%\), BS-AG about \(49\)--\(64\%\) of the canonical
headroom at \(B\in\{2,3,4\}\). &
Answered positively. \\
\addlinespace
Are interactions present? &
Yes.  Pair interactions are structured and mostly negative / redundant.
Cheap pairwise composition helps at \(B=2\) but fails at \(B=3,4\). &
Answered as diagnostic result. \\
\addlinespace
Is the landscape smooth enough for BO or submodular theory? &
Not clearly.  One-swap smoothness exists but is weak, diminishing-return gaps
are mildly positive but region-dependent, and BO remains unclear. &
Answered conservatively. \\
\bottomrule
\end{tabular}
\end{center}

\noindent
This means the current thesis is not a speculative proposal from zero.  It is
closer to a synthesis and formalization of an empirical story already visible
in the results:

\begin{quote}
Correction timing matters on ProSeCo-OWT.  Simple observable uncertainty
signals do not reliably predict the best timing.  The failure is not random:
there is structured redundancy and local search structure, but not enough
evidence for a clean submodular or Bayesian-optimization theory.
\end{quote}

\section*{Model choice and external validity}

\begin{center}
\small
\begin{tabular}{p{0.22\linewidth}p{0.18\linewidth}p{0.47\linewidth}}
\toprule
Model / codebase & Role & Reason \\
\midrule
ProSeCo-OWT & Main case study &
Best fit.  It is a self-correcting masked-diffusion backend with an available
small OpenWebText checkpoint and a fixed correction mechanism.  It lets the
thesis study timing without training a corrector from scratch. \\
\addlinespace
MDLM-OWT & Baseline only &
Feasible and instrumentable, but only relevant if a meaningful correction
intervention is defined.  Without that, it is a predictor model rather than a
self-correction timing platform. \\
\addlinespace
Informed-correctors code & Controlled fallback &
Theoretically aligned with predictor-corrector sampling and informed
correctors.  Less aligned with language-scale ProSeCo, so best as a small
controlled appendix if needed. \\
\addlinespace
SEDD & Optional comparison &
Important discrete diffusion model, but not naturally a ProSeCo-style
self-correcting masked sampler. \\
\addlinespace
ProSeCo-LLaDA / LLaDA 8B & Stretch only &
Scientifically relevant but infrastructure-heavy.  Existing LLaDA evidence in
this project did not show clear positive oracle headroom, so it should not be
the core model. \\
\bottomrule
\end{tabular}
\end{center}

\paragraph{Does it make sense to try the heavier model?}
Only as an external-validity gate, not as the main path.  The heavier
LLaDA/ProSeCo-LLaDA direction is scientifically attractive, but it has high
infrastructure risk and current local evidence is inconclusive: the bounded
LLaDA-SFT probe did not transfer the positive MC-oracle headroom at tested
resolution.  Before running a full heavy-model scheduling study, the necessary
gate is a small pilot showing:
\begin{itemize}
\item positive correction headroom over uniform with a meaningful confidence
interval;
\item non-degenerate corrector behavior;
\item a stable and interpretable utility metric.
\end{itemize}
Without this gate, a heavier model is likely to consume time without improving
the thesis argument.

\section*{What is model-agnostic, and what is not}

The theory is model-agnostic in a precise but limited sense.  The definitions
of \(J(S)\), \(G(S)\), \(\Delta_t\), \(m(t,z)\), additivity error, pair
interactions, and finite-pool regret apply to any fixed tuple:
\[
\text{(predictor, corrector, utility }F\text{, budget }B\text{)}.
\]
They do not require ProSeCo, reversibility, detailed balance, or a specific
neural architecture.

The empirical classification is not model-agnostic.  The statement
``simple signals fail and search recovers headroom'' is currently a
ProSeCo-OWT result.  A different model/corrector/metric could lie in a
different regime:
\begin{itemize}
\item no-op: correction has no measurable headroom;
\item marginal: time or uncertainty signals predict correction value well;
\item interaction/redundancy: marginal signals are insufficient but
interactions are structured;
\item higher-order/search: only direct true-\(G\) search finds the good
schedules.
\end{itemize}

\noindent
Thus the model-agnostic output of the thesis is the diagnostic framework and
the finite-pool certificates.  The ProSeCo-OWT output is a case-study
classification within that framework.

\section*{Recommended next steps}

Given the work already completed, the evidence does not support starting a
major new online scheduler project.  The simple online/state conditioning and
held-out state-predictability tests are negative, and the landscape evidence
says Bayesian optimization is still unclear.  The most defensible next steps
are:

\begin{enumerate}
\item \textbf{Write the thesis story around the results already obtained.}
The current results already answer the main diagnostic questions on
ProSeCo-OWT.  The immediate task is synthesis, not another large experiment.
\item \textbf{Do not build an online scheduler from the current feature set.}
An online scheduler is justified only if current-state features add held-out
predictive information beyond time.  The completed audit does not show this.
\item \textbf{Use the state-predictability negative as a result.}
The negative audit is useful evidence: it shows that the current observable
uncertainty/error features do not explain correction timing enough to support
a simple adaptive trigger.
\item \textbf{Use heavier models only for a feasibility gate.}
Run LLaDA/ProSeCo-LLaDA only if external validity becomes mandatory and only
after a small headroom pilot passes.
\item \textbf{Prepare the supervisor discussion around scope.}
The key decision is whether a ProSeCo-OWT case study plus model-agnostic
diagnostic framework is accepted as the thesis contribution, or whether the
supervisor requires one additional external-validity check.
\end{enumerate}

\section*{Source anchors}

\begin{itemize}
\item Zhao, Shi, Chen, Druckmann, Mackey, and Linderman,
\emph{Informed Correctors for Discrete Diffusion Models},
\href{https://arxiv.org/abs/2407.21243}{arXiv:2407.21243}.
\item Zanella, \emph{Informed proposals for local MCMC in discrete spaces},
\href{https://arxiv.org/abs/1711.07424}{arXiv:1711.07424} / JASA 2020.
\item Lavenant and Zanella,
\emph{Error Bounds and Optimal Schedules for Masked Diffusions with Factorized
Approximations},
\href{https://arxiv.org/abs/2510.25544}{arXiv:2510.25544}.
\item Kim et al.,
\emph{Fine-Tuning Masked Diffusion for Provable Self-Correction},
\href{https://arxiv.org/abs/2510.01384}{arXiv:2510.01384}.
\item Schiff et al.,\\
\emph{Learn from Your Mistakes: Self-Correcting Masked Diffusion Models},\\
\href{https://arxiv.org/abs/2602.11590}{arXiv:2602.11590}.
\end{itemize}

\end{document}
